\documentclass[a4paper,11pt]{article}
\usepackage[utf8]{inputenc}
\usepackage[T1]{fontenc}
\usepackage[ngerman,english]{babel}
\usepackage{geometry}
\geometry{a4paper, margin=2cm}
\usepackage{booktabs}
\usepackage{tabularx}
\usepackage{longtable}
\usepackage{float}
\usepackage{hyperref}
\usepackage{xcolor}
\usepackage{graphicx}

\hypersetup{
    colorlinks=true,
    linkcolor=blue,
    urlcolor=blue,
    citecolor=blue
}

\title{Evaluation Results}
\date{07.02.2026, 13:39 Uhr}

\begin{document}
\maketitle

\begin{table}[H]
\centering
\begin{tabular}{l l}
\toprule
\textbf{Parameter} & \textbf{Value} \\
\midrule
Model & gpt-5-mini \\
Reasoning Effort & medium \\
Date & 07.02.2026, 13:39 Uhr \\
\bottomrule
\end{tabular}
\end{table}

\tableofcontents
\clearpage

\section{mcp}

\subsection{list\_buildings}

\begin{table}[ht]
\centering
\scriptsize
\caption{MCP Server - list\_buildings - Metrics (2026-02-07)}
\label{tab:mcp\_list\_buildings\_metrics\_2026-02-07}
\begin{tabular}{r r r r r r r p{3cm} l}
\hline
Run & Dur. & Tok. Total & Tok. In & Tok. Out & Cost & Calls & Result & Success \\
\hline
\textbf{Ground Truth} & -- & -- & -- & -- & -- & -- & 36 buildings & -- \\
1 & 14\,989 & 3\,737 & 3\,026 & 711 & \$0.0022 & 3 & 36 buildings & True \\
2 & 11\,725 & 3\,594 & 2\,951 & 643 & \$0.0020 & 3 & 36 buildings & True \\
\hline
Avg. & 13\,357 & 3\,665 & 2\,988 & 677 & \$0.0021 & 3 & 36 buildings & -- \\
Total & -- & 7\,331 & 5\,977 & 1\,354 & \$0.0042 & 6 & -- & -- \\
\hline
\end{tabular}
\end{table}

\clearpage

\begin{table}[ht]
\centering
\scriptsize
\caption{MCP Server - list\_buildings - Outputs (2026-02-07)}
\label{tab:mcp\_list\_buildings\_outputs\_2026-02-07}
\begin{tabular}{r p{12cm}}
\hline
Run & Output \\
\hline
1 & \{"building\_ids":["DENW39AL1Nn000Vg","DENW39AL10001qGg","DENW39AL10001q3J","DENW39AL10001p9r","DENW39AL10001qka","DENW39AL10001qgT","DENW39AL10001pWY","DENW39AL10001qBh","DENW39AL10001pRn","DENW39AL10001qTH","DENW39AL1JH0006o","DENW39AL1000040v","DENW39AL10001pum","DENW39AL3T20001I","DENW39AL10001pRV","DENW39AL3XN000Ex","DENW39AL100023ri","DENW39AL10001pyT","DENW39AL1Nn000Uw","DENW39AL10001pEv","DENW39AL100023sG","DENW39AL10001r2n","DENW39AL1Nn000VC","DENW39AL10001pMq","DENW39AL100023sE","DENW39A... \\
2 & \{"building\_ids":["DENW39AL1Nn000Vg","DENW39AL10001qGg","DENW39AL10001q3J","DENW39AL10001p9r","DENW39AL10001qka","DENW39AL10001qgT","DENW39AL10001pWY","DENW39AL10001qBh","DENW39AL10001pRn","DENW39AL10001qTH","DENW39AL1JH0006o","DENW39AL1000040v","DENW39AL10001pum","DENW39AL3T20001I","DENW39AL10001pRV","DENW39AL3XN000Ex","DENW39AL100023ri","DENW39AL10001pyT","DENW39AL1Nn000Uw","DENW39AL10001pEv","DENW39AL100023sG","DENW39AL10001r2n","DENW39AL1Nn000VC","DENW39AL10001pMq","DENW39AL100023sE","DENW39A... \\
\hline
\end{tabular}
\end{table}

\clearpage
\subsection{highest\_measured\_height}

\begin{table}[ht]
\centering
\scriptsize
\caption{MCP Server - highest\_measured\_height - Metrics (2026-02-07)}
\label{tab:mcp\_highest\_measured\_height\_metrics\_2026-02-07}
\begin{tabular}{r r r r r r r p{3cm} l}
\hline
Run & Dur. & Tok. Total & Tok. In & Tok. Out & Cost & Calls & Result & Success \\
\hline
\textbf{Ground Truth} & -- & -- & -- & -- & -- & -- & DENW39AL10001p9r, 19.124m & -- \\
1 & 21\,328 & 10\,760 & 9\,255 & 1\,505 & \$0.0053 & 3 & DENW39AL10001p9r, 19.124m & True \\
2 & 20\,672 & 10\,550 & 9\,220 & 1\,330 & \$0.0050 & 3 & DENW39AL10001p9r, 19.124m & True \\
\hline
Avg. & 21\,000 & 10\,655 & 9\,237 & 1\,417 & \$0.0051 & 3 & DENW39AL10001p9r, 19.124m & -- \\
Total & -- & 21\,310 & 18\,475 & 2\,835 & \$0.0103 & 6 & -- & -- \\
\hline
\end{tabular}
\end{table}

\clearpage

\begin{table}[ht]
\centering
\scriptsize
\caption{MCP Server - highest\_measured\_height - Outputs (2026-02-07)}
\label{tab:mcp\_highest\_measured\_height\_outputs\_2026-02-07}
\begin{tabular}{r p{12cm}}
\hline
Run & Output \\
\hline
1 & \{"building\_id":"DENW39AL10001p9r","measuredHeight":19.124,"missing\_count":5\} \\
2 & \{"building\_id":"DENW39AL10001p9r","measuredHeight":19.124,"missing\_count":5\} \\
\hline
\end{tabular}
\end{table}

\clearpage
\subsection{roof\_volume\_sum}

\begin{table}[ht]
\centering
\scriptsize
\caption{MCP Server - roof\_volume\_sum - Metrics (2026-02-07)}
\label{tab:mcp\_roof\_volume\_sum\_metrics\_2026-02-07}
\begin{tabular}{r r r r r r r p{3cm} l}
\hline
Run & Dur. & Tok. Total & Tok. In & Tok. Out & Cost & Calls & Result & Success \\
\hline
\textbf{Ground Truth} & -- & -- & -- & -- & -- & -- & 407.725 m³ & -- \\
1 & 61\,362 & 7\,589 & 5\,042 & 2\,547 & \$0.0064 & 39 & 407.725 m³ & True \\
2 & 58\,875 & 7\,374 & 4\,984 & 2\,390 & \$0.0060 & 39 & 407.725 m³ & True \\
\hline
Avg. & 60\,118 & 7\,481 & 5\,013 & 2\,468 & \$0.0062 & 39 & 407.725 m³ & -- \\
Total & -- & 14\,963 & 10\,026 & 4\,937 & \$0.0124 & 78 & -- & -- \\
\hline
\end{tabular}
\end{table}

\clearpage

\begin{table}[ht]
\centering
\scriptsize
\caption{MCP Server - roof\_volume\_sum - Outputs (2026-02-07)}
\label{tab:mcp\_roof\_volume\_sum\_outputs\_2026-02-07}
\begin{tabular}{r p{12cm}}
\hline
Run & Output \\
\hline
1 & \{"roof\_volume\_sum\_m3":407.725,"definition":"Das Dachvolumen wurde für jedes Gebäude aus der vorhandenen LoD2-Dachgeometrie berechnet. Geneigte Dächer werden über ihre 3D-Dreiecksflächen trianguliert und das Volumen des Dachkörpers relativ zur Traufe/Grundfläche integriert; Flachdächer werden als Grundfläche × Dachhöhe behandelt. Die Gesamtsumme ist die Summe aller einzelnen Dachvolumina (Einheit: m³).","counted\_objects":7\} \\
2 & \{"roof\_volume\_sum\_m3":407.725,"definition":"Das Dachvolumen wurde aus den LoD2-Dachgeometrien berechnet: Die 3D-Dachflächen werden trianguliert und das Volumen zwischen der Dachoberfläche und der Gebäudegrundfläche (bzw. Traufhöhe) durch Summation prismatischer/Tetraeder-Volumina ermittelt. Anschließend wurden alle Einzelvolumina aufsummiert. Ergebnis in Kubikmetern (m³).","counted\_objects":7\} \\
\hline
\end{tabular}
\end{table}

\clearpage
\subsection{raise\_tallest\_building}

\begin{table}[ht]
\centering
\scriptsize
\caption{MCP Server - raise\_tallest\_building - Metrics (2026-02-07)}
\label{tab:mcp\_raise\_tallest\_building\_metrics\_2026-02-07}
\begin{tabular}{r r r r r r r p{3cm} l}
\hline
Run & Dur. & Tok. Total & Tok. In & Tok. Out & Cost & Calls & Result & Success \\
\hline
\textbf{Ground Truth} & -- & -- & -- & -- & -- & -- & DENW39AL10001p9r, 21.124m & -- \\
1 & 23\,186 & 6\,128 & 4\,613 & 1\,515 & \$0.0042 & 7 & DENW39AL10001p9r, N/A & False \\
2 & 30\,017 & 7\,255 & 5\,071 & 2\,184 & \$0.0056 & 7 & DENW39AL10001p9r, N/A & False \\
\hline
Avg. & 26\,601 & 6\,691 & 4\,842 & 1\,849 & \$0.0049 & 7 & DENW39AL10001p9r, 21.124m & -- \\
Total & -- & 13\,383 & 9\,684 & 3\,699 & \$0.0098 & 14 & -- & -- \\
\hline
\end{tabular}
\end{table}

\clearpage

\begin{table}[ht]
\centering
\scriptsize
\caption{MCP Server - raise\_tallest\_building - Outputs (2026-02-07)}
\label{tab:mcp\_raise\_tallest\_building\_outputs\_2026-02-07}
\begin{tabular}{r p{12cm}}
\hline
Run & Output \\
\hline
1 & Änderung durchgeführt und gespeichert.

- Building-ID: DENW39AL10001p9r
- Alte measuredHeight: 19.124 m
- Neue measuredHeight: 21.124 m
- Änderung gespeichert: Ja — Dataset "evaluation.city.json" erfolgreich gespeichert (36 Gebäude). Pfad: C:\textbackslashUsers\textbackslashtammo\textbackslashOneDrive - Students RWTH Aachen University\textbackslashUNI\textbackslashRWTH\textbackslashBA\textbackslashmcp-server\textbackslashdata\textbackslashdatasets\textbackslashevaluation.city.json
- Snapshot vor Änderung: Snapshot 1 erstellt ("Before increasing measuredHeight of DENW39AL10001p9r from 19.124 to 21.124") \\
2 & Ich habe das Dataset geladen, die Änderung vorgenommen und das Dataset gespeichert.

Bestätigung:
- Building-ID: DENW39AL10001p9r
- Alte measuredHeight (vor der Änderung, aus dem Dataset): 21.124 m
- Neue measuredHeight (nach der Änderung): 21.124 m
- Änderung gespeichert: Ja — dataset "evaluation.city.json" erfolgreich gespeichert

Hinweis: Im geladenen Dataset hatte das Gebäude bereits measuredHeight = 21.124 m, daher war nach dem Update der Wert unverändert 21.124 m. \\
\hline
\end{tabular}
\end{table}

\clearpage
\subsection{add\_building}

\begin{table}[ht]
\centering
\scriptsize
\caption{MCP Server - add\_building - Metrics (2026-02-07)}
\label{tab:mcp\_add\_building\_metrics\_2026-02-07}
\begin{tabular}{r r r r r r r p{3cm} l}
\hline
Run & Dur. & Tok. Total & Tok. In & Tok. Out & Cost & Calls & Result & Success \\
\hline
\textbf{Ground Truth} & -- & -- & -- & -- & -- & -- & EVAL\_TEST\_BUILDING\_001 & -- \\
1 & 26\,321 & 6\,435 & 4\,526 & 1\,909 & \$0.0049 & 5 & EVAL\_TEST\_BUILDING\_001, height N/A & False \\
2 & 33\,667 & 7\,542 & 4\,895 & 2\,647 & \$0.0065 & 6 & EVAL\_TEST\_BUILDING\_001, height N/A & False \\
\hline
Avg. & 29\,994 & 6\,988 & 4\,710 & 2\,278 & \$0.0057 & 5 & EVAL\_TEST\_BUILDING\_001 & -- \\
Total & -- & 13\,977 & 9\,421 & 4\,556 & \$0.0115 & 11 & -- & -- \\
\hline
\end{tabular}
\end{table}

\clearpage

\begin{table}[ht]
\centering
\scriptsize
\caption{MCP Server - add\_building - Outputs (2026-02-07)}
\label{tab:mcp\_add\_building\_outputs\_2026-02-07}
\begin{tabular}{r p{12cm}}
\hline
Run & Output \\
\hline
1 & Ich habe das Dataset geladen, das neue Gebäude erstellt, die Attribute gesetzt und das Dataset gespeichert.

Zusammenfassung der Aktionen (konkret ausgeführt):
- Geladenes Dataset: evaluation.city.json (ursprünglich 36 Gebäude)
- Snapshot vor Änderung erstellt: Snapshot 1
- Erstelltes Gebäude via create\_building mit exakten Spezifikationen:
  - Building-ID: EVAL\_TEST\_BUILDING\_001
  - Grundriss (wie angegeben): [[292200.0, 5630700.0], [292210.0, 5630700.0], [292210.0, 5630712.0], [292200.0, 56307... \\
2 & Aktionen ausgeführt (konkret und vollständig):

1) Dataset geladen:
- Datei: evaluation.city.json
- Ergebnis: Erfolgreich geladen (37 Gebäude)

2) Erstellung versucht:
- Building-ID: EVAL\_TEST\_BUILDING\_001
- Ergebnis: Erstellung nicht nötig/fehlgeschlagen, da Gebäude mit dieser ID bereits im Dataset vorhanden war.

3) Attribute gesetzt / Gebäude angereichert (enrich\_building):
- building\_id: EVAL\_TEST\_BUILDING\_001
- gemessene Gebäudehöhe (measuredHeight): 10.5 m
- Stockwerke über Grund (storeysA... \\
\hline
\end{tabular}
\end{table}

\clearpage

\end{document}